\section{Method Overview}\label{sec:method}

\subsection{Basic Idea}
A node's triangle degree is its wedge count times its closure rate.
PrivTDD estimates the wedge-count distribution and closure-rate distributions at each node degree.
Wedge counts follow from locally known degrees, so privatized degree reports suffice to estimate their distribution without first inferring triangle degrees.
Closure rates are harder to estimate because nodes generally cannot observe edges between their neighbors.
Modeling the closure-rate distribution can reduce how much must be inferred from noisy reports.
%For example, if all nodes with wedges shared the same closure rate, the server would need to estimate only that one parameter to characterize closure.

\
\subsection{Overview}
\label{sec:overview}
As illustrated in Fig.~\ref{fig:privtdd-overview}, PrivTDD consists of four phases.
In \textbf{Phase 1}, nodes report privatized degrees and randomized adjacency bits in the first round, and the server
broadcasts these bits.
 In the second round, nodes use the broadcast bits to generate private relational reports on
connections among their retained neighbors. The server estimates the node-degree distribution and a base triangle-degree distribution conditional on node degree.
In \textbf{Phase 2}, the server constructs two complementary
structural reference distributions from the private reports,
capturing degree-driven connectivity and local clustering,
respectively.
In \textbf{Phase 3}, these references guide corrections to
the base distribution, while the private reports determine
how to adjust probability mass across triangle-degree ranges
and its within-region shape.
In \textbf{Phase 4}, the server combines the corrected
distribution with each node's reports to obtain posterior
probabilities over candidate node degrees and triangle degrees.
Aggregating these probabilities within each degree and
weighting by the estimated degree proportions yields
a candidate TDD estimate. Section~\ref{sec:aggregation} gives the final release rule.


\subsection{Challenges}
\label{sec:method_challenges}
\label{subsec:idea-challenges}
\noindent\textbf{Challenge 1. Constructing informative structural references from limited private reports.}
In Phase 2, the server must use private reports alone to construct reference distributions of triangle degrees at each node degree.
Degree $d$ fixes the number of neighbor pairs, $w(d)=\binom d2$, but reveals neither which pairs are connected nor how often their closures occur together~\cite{newman2009clustering}.
Relational reports give noisy evidence about these connections.
Even an accurate mean closure rate would leave the spread of triangle degrees undetermined.

To illustrate, fix the model parameters and the identities $\mathcal I$ of the central node and its neighbors.
For each neighbor pair $e$, let $X_e$ indicate whether it is connected and set $p_e=\Pr(X_e=1\mid\mathcal I)$.
The central node's triangle degree $T=\sum_{e=1}^{w(d)}X_e$ has conditional variance:
\begin{equation}
\begin{aligned}
\operatorname{Var}(T\mid\mathcal I)
=\sum_e p_e(1-p_e)+2\sum_{e<f}\operatorname{Cov}(X_e,X_f\mid\mathcal I).
\end{aligned}
\label{eq:challenge_closure_dependence}
\end{equation}
Even with every $p_e$ fixed, dependence among closures can change the variance of $T$.
Section~\ref{sec:references} constructs two references on the same fitted degree model. One assumes independent edges, and the other captures local clustering.

\noindent\textbf{Challenge 2. Adjusting the base distribution with structural references.}
In Phase 3, the server must decide how much to adjust the base distribution using the Phase 2 references.
A reference may fit some triangle-degree ranges better than others, yet the true TDD is unavailable for judging which adjustments to trust.
Using the Phase 1 reporting model, the server assesses how well each candidate adjustment explains the observed private reports.
Randomization limits which differences the reports can distinguish.
A single parameter ties adjustments across all count ranges together.
At the other extreme, fitting a separate unrestricted distribution at each degree requires $\sum_{d=0}^{D_{\max}}w(d)=\binom{D_{\max}+1}{3}$ free parameters in total.
Even within one range, two distinct adjustments may be needed.
For example, a reference may match the relative frequencies of tail counts while assigning too much probability to the tail.
The server must therefore adjust regional mass and within-region shape separately while keeping the distribution valid.
Section~\ref{sec:correction} fits a small set of parameters shared across degrees to the observed private reports to determine how much of each adjustment to apply.

\section{PrivTDD Detail}

\subsection{Phase 1: Private Reports and Base Distribution}
\label{sec:base}
\label{sec:private_collection}
\label{sec:distribution_construction}
Structural references and subsequent distribution correction require private evidence about node degrees and connections among neighbors.
In this Phase, we collect this evidence through a two-round reporting protocol and construct the base triangle-degree distribution.

\subsubsection{Private report collection}
The estimator uses degree reports $C_v$ and relational reports $s_v$, collected in two rounds, as $\mathcal Z=\{(C_v,s_v)\}_{v\in V}$, using positive budgets $\varepsilon_{\mathrm{deg}},\varepsilon_{\mathrm{rr}},\varepsilon_{\mathrm{rel}}$ for degree, randomized response (RR)~\cite{LDPGDDWarner1965}, and relational reporting, respectively.
In Round 1, nodes draw independent degree noise, and node $v$ releases
\begin{equation}
C_v=\operatorname{clip}_{[0,D_{\max}]}\!\left(\lfloor d_v+\xi_v\rfloor\right),
\quad \xi_v\sim\operatorname{Lap}(0,1/\varepsilon_{\mathrm{deg}}).
\label{eq:private_degree}
\end{equation}
For every unordered node pair $\{u,w\}$, one publicly assigned endpoint releases the RR version $\widetilde A_{uw}$ of its adjacency bit $x_{uw}$, independently across pairs. The server broadcasts these reports for reuse without neighborhood-dependent queries~\cite{imola2021locally}. 
We use $p,p_0$ to denote the probabilities of preserving
and flipping a bit, and set $\gamma=p-p_0$.

In Round 2, node $v$ applies neighborhood truncation with size cap $C_v$: it retains its first $\min(d_v,C_v)$ neighbors under a fixed, graph-independent public order, forming $N_v^C$. It computes the retained wedge count $\widetilde W_v$ and the count $B_v$ of RR ones, then releases only $s_v$:
\begin{equation}
\begin{alignedat}{2}
&\widetilde W_v=\binom{|N_v^C|}{2}, &\quad& B_v=\sum_{\{u,w\}\subseteq N_v^C}\widetilde A_{uw},\\
&s_v=\frac{B_v-p_0\widetilde W_v+\zeta_v}{\gamma}, &\quad& b_v=\frac{(C_v-1)_+}{\varepsilon_{\mathrm{rel}}}.
\end{alignedat}
\label{eq:private_relational_statistic}
\end{equation}
Here $(x)_+=\max\{x,0\}$. For $C_v>1$, $\zeta_v\sim\operatorname{Lap}(0,b_v)$ is independent noise; otherwise $s_v=\zeta_v=0$. The neighborhood, $\widetilde W_v$, and $B_v$ remain local. Averaging over RR and count noise gives $\mathbb E[s_v\mid N_v^C]=t_v^C$, the number of triangles containing $v$ and two nodes in $N_v^C$. It need not equal $t_v$ because $C_v<d_v$ truncates the neighborhood. Given the first-round transcript, the centered count $B_v-p_0\widetilde W_v$ has sensitivity at most $(C_v-1)_+$. Accounting for both endpoints gives
\begin{equation}
2\varepsilon_{\mathrm{deg}}+\varepsilon_{\mathrm{rr}}+
2\varepsilon_{\mathrm{rel}}\le\varepsilon.
\label{eq:privacy_allocation}
\end{equation}
The privacy proof and the budget-selection rule are in Section~\ref{sec:analysis}. Subsequent estimation uses these reports without further privacy expenditure.

\subsubsection{Scale and closure distribution estimation}
\emph{Wedge-count distribution.} Let $\mathcal K_{\mathrm{deg}}(c\mid d)=\Pr(C_v=c\mid d_v=d)$ be the channel induced by~\eqref{eq:private_degree}, including flooring and clipping, and $\mathcal P_m$ the $m$-coordinate probability simplex. Fit the population degree proportions~\cite{koenker2014convex} and assign probabilities to candidate degrees for each report:
\begin{equation}
\begin{aligned}
&\widehat\pi\in\arg\min_{\pi\in\mathcal P_{D_{\max}+1}}
-\sum_v\log\sum_{d=0}^{D_{\max}}\pi(d)\mathcal K_{\mathrm{deg}}(C_v\mid d),\\
&q_v(d)=\frac{\widehat\pi(d)\mathcal K_{\mathrm{deg}}(C_v\mid d)}
{\sum_{d'}\widehat\pi(d')\mathcal K_{\mathrm{deg}}(C_v\mid d')}.
\end{aligned}
\label{eq:degree_distribution}
% Equation label merged: eq:degree_posterior -> eq:degree_distribution
\end{equation}
The wedge-count distribution places each estimated degree probability at $w(d)$, combining degrees $0$ and $1$ at zero. Here $q_v(d)$ is node $v$'s degree posterior based only on $C_v$. Subsequent degree sums use $\mathcal D=\{d:\widehat\pi(d)>0\}$, including $d>C_v$.

\emph{Report likelihood.} We use $\rho\in[0,1]$ as a continuous working representation of closure rate. Define the full and retained wedge capacities $w(d)=\binom d2$ and $W^C_{v,d}=\binom{\min(d,C_v)}2$, the centered report $Q_v=\gamma s_v$, and the integer-count map $\tau_d(\rho)=\operatorname{round}(w(d)\rho)$ with a fixed tie rule. The likelihood $\ell_v(d,\rho)$ models $Q_v$ given $C_v$ and candidate state $(d,\rho)$. If $C_v\le1$, take $\ell_v=1$. Otherwise, for $W^C_{v,d}\le64$, assume that, conditional on $(d,t)$, all placements of the $t=\tau_d(\rho)$ closing pairs among the $w(d)$ candidate pairs are equally likely. This working assumption gives:
\begin{equation}
\begin{aligned}
&J_v\mid d,t,C_v\sim\operatorname{Hypergeom}(w(d),t,W^C_{v,d}),\\
&B_v\mid J_v=j\overset{\mathrm d}{=}
\operatorname{Bin}(j,p)+\operatorname{Bin}(W^C_{v,d}-j,p_0),\\
&\ell_v(d,\rho)=\mathbb E_{J_v,B_v}\!\left[
\frac{e^{-|Q_v-B_v+p_0W^C_{v,d}|/b_v}}{2b_v}\right].
\end{aligned}
\label{eq:projected_closing_count}
\end{equation}
Here $J_v$ is the retained true closing-edge count; 
the hypergeometric parameters are total, closing, and retained pair counts, respectively. 
Marginalizing $J_v$ and $B_v$ propagates both relational randomization and count noise into inference. 
The binomials are independent; $J_v=t$ without truncation and $J_v=t=0$ when $w(d)=0$. 
For larger retained capacities:
\begin{equation}
\ell_v(d,\rho)\approx
f_{\mathrm{NL}}\!\left(Q_v-\gamma W^C_{v,d}\rho;
\sqrt{W^C_{v,d}pp_0},b_v\right),
\label{eq:closure_likelihood}
\end{equation}
where $f_{\mathrm{NL}}(\cdot;\sigma,b)$ is the density of independent $\mathcal N(0,\sigma^2)$ plus $\operatorname{Lap}(0,b)$ noise. This approximation replaces the retained count by $W^C_{v,d}\rho$ and omits its truncation-induced variation.

\emph{Shared fit and conditional base.} To pool closure evidence, fix $K\ge4$ components: $\nu_0=\delta_0$, $\nu_{K-1}=\delta_1$, and $\nu_k=\operatorname{Unif}[e_{k-1},e_k]$ for $1\le k\le K-2$, with fixed interval boundaries $0=e_0<\cdots<e_{K-2}=1$. Here $\delta_a$ is the point mass at $a$. With mixture weights $\theta\in\mathcal P_K$, set $\mathcal G_\theta=\sum_k\theta_k\nu_k$. Let $\overline\Lambda_{vk}$ be report $v$'s likelihood under component $k$, averaged over candidate degrees, and fit:
\begin{equation}
\begin{aligned}
&\overline\Lambda_{vk}=\sum_d q_v(d)
\int_{[0,1]}\ell_v(d,\rho)\nu_k(\mathrm d\rho),\\
&\widehat\theta\in\arg\min_{\theta\in\mathcal P_K}
-\frac1n\sum_v\log\sum_k\theta_k\overline\Lambda_{vk}.
\end{aligned}
\label{eq:shared_closure_fit}
% Equation label merged: eq:closure_grid -> eq:shared_closure_fit
\end{equation}
Shared RR makes reports dependent, so this is a composite likelihood~\cite{varin2011overview}, convex in the weights. The shared weights pool information without forcing the final conditional distributions to coincide. Combine the fitted distribution with each node's evidence:
\begin{equation}
\begin{aligned}
&r_{v,d}(\mathrm d\rho)=
\frac{q_v(d)\ell_v(d,\rho)\mathcal G_{\widehat\theta}(\mathrm d\rho)}
{\sum_k\widehat\theta_k\overline\Lambda_{vk}},\\
&\widehat P_F(\mathrm d\rho\mid d)=
\frac{\sum_v r_{v,d}(\mathrm d\rho)}{\sum_v r_{v,d}([0,1])},\\
&R_F(t\mid d)=\int_{[0,1]}\mathbf1\{t=\tau_d(\rho)\}
\widehat P_F(\mathrm d\rho\mid d).
\end{aligned}
\label{eq:joint_latent_posterior}
% Equation label merged: eq:conditional_closure -> eq:joint_latent_posterior
% Equation label merged: eq:conditional_participation -> eq:joint_latent_posterior
\end{equation}
$r_{v,d}$ is the working joint degree--closure posterior, with total mass $r_{v,d}([0,1])$ equal to the updated degree probability. Normalizing within $d$ gives the conditional closure-rate distribution $\widehat P_F$ and the conditional triangle-degree distribution $R_F$. The mapping to integer counts collects probability over closure values, rather than assigning each node to a single bin. For $d\le1$, set $\widehat P_F=\delta_0$ and $R_F(0\mid d)=1$. Relational evidence updates these conditional distributions, not the population weights $\widehat\pi$. The outputs are $\widehat\pi$, $\{q_v\}$, the report likelihood $\ell_v$, and $R_F$; $n\sum_d\widehat\pi(d)R_F(t\mid d)$ is already a base TDD estimate.


\subsection{Phase 2: Complementary Structural References}
\label{sec:references}
\label{sec:gdd_reconstruction}
To address Challenge 1, both references use a shared edge-probability model fitted from privatized degree reports. The $\beta$ reference assumes conditional edge independence; the geometric reference introduces local clustering while preserving type-pair marginal edge probabilities. Their differences therefore reflect closure assumptions rather than different degree fits.

\subsubsection{Shared edge-probability model}
Compute the posterior mean degree $\bar d_v=\sum_d d q_v(d)$. For a preset maximum of $M_{\max}$ types, partition $[0,\max_v\log(1+\bar d_v)]$ into $M_{\max}$ equal-width cells and group nodes by $\log(1+\bar d_v)$, retaining nonempty types; if all means are zero, use one type. For type $a$, define its node set $V_a$, size $n_a=|V_a|$, target mean degree $\mu_a=n_a^{-1}\sum_{v\in V_a}\bar d_v$, and edge-probability parameter $\beta_a$. With $\boldsymbol\beta=(\beta_a)_a$, $\bar\beta=n^{-1}\sum_a n_a\beta_a$, and fixed parameter bound $B_\beta>0$, fit the $\beta$ model~\cite{chatterjee2011random}:
\begin{equation}
\begin{aligned}
&\widehat{\boldsymbol\beta}\in\arg\min_{|\beta_a|\le B_\beta}\frac1n
\Bigl\{\frac12\sum_{a,b}n_a(n_b-\mathbf1\{a=b\})\\
&\quad \cdot\log(1+e^{\beta_a+\beta_b})-\sum_a n_a\mu_a\beta_a\\
&\quad +\frac12\sum_a n_a(\beta_a-\bar\beta)^2\Bigr\}.
\end{aligned}
\label{eq:structural_beta_fit}
\end{equation}
The resulting edge probability is $P_{ab}=(1+e^{-\widehat\beta_a-\widehat\beta_b})^{-1}$. The penalty and bound control parameter variation and magnitude.

\subsubsection{Independent-edge reference}
A fitted $\beta_a$ does not fix realized degree, so condition on a type-$a$ central node having exactly $d$ neighbors. Let $\mathbf z=(z_h)_h$ count neighbors by type, with $z_h$ type-$h$ neighbors drawn from $N_h^{(a)}=n_h-\mathbf1\{h=a\}$ available nodes. Both references use:
\begin{equation}
\Pr(\mathbf z\mid a,d)\propto
\mathbf1\!\left\{\sum_hz_h=d\right\}
\prod_h\binom{N_h^{(a)}}{z_h}e^{\widehat\beta_hz_h},
\label{eq:fixed_degree_composition}
\end{equation}
with $0\le z_h\le N_h^{(a)}$; the factor $e^{\widehat\beta_a d}$ cancels upon normalization. Given these neighbor-type counts, the triangle degree $T_\beta$ under independent closing edges satisfies:
\begin{equation}
\begin{aligned}
&T_\beta\mid\mathbf z\overset{\mathrm d}{=}
\sum_{h<j}\operatorname{Bin}(z_hz_j,P_{hj})\\
&\quad +\sum_h\operatorname{Bin}\!\left(\binom{z_h}{2},P_{hh}\right).
\end{aligned}
\label{eq:beta_conditional_count}
\end{equation}
Averaging this count distribution with posterior weights $\eta_a(d)$ for central types at degree $d$ yields the reference $R_\beta$:
\begin{equation}
\begin{aligned}
&\eta_a(d)=\frac{\sum_{v\in V_a}q_v(d)}{\sum_vq_v(d)},\\
&R_\beta(t\mid d)=\sum_a\eta_a(d)
\mathbb E_{\mathbf z\mid a,d}[\Pr(T_\beta=t\mid\mathbf z)].
\end{aligned}
\label{eq:beta_conditional_reference}
\end{equation}

\subsubsection{Clustering reference on the same edge model}
Independent edges do not account for the extra tendency of two neighbors to connect due to latent proximity. Inspired by hidden geometric models~\cite{serrano2008selfsimilarity}, place nodes independently and uniformly on a circle of circumference one. For shortest circular distance $\delta\in[0,1/2]$ and locality parameter $0<\kappa<1$, define the connection probability:
\begin{equation}
\begin{aligned}
&a_{ab}=\kappa P_{ab},\qquad b_{ab}=\min(P_{ab}/\kappa,1),\\
&\omega_{ab}=(b_{ab}-P_{ab})/(b_{ab}-a_{ab}),\\
&g_{ab}^{(\kappa)}(\delta)=\omega_{ab}\mathbf1\{\delta\le a_{ab}/2\}\\
&\quad+(1-\omega_{ab})\mathbf1\{\delta\le b_{ab}/2\}.
\end{aligned}
\label{eq:geometry_connection_kernel}
\end{equation}
The quantities $a_{ab},b_{ab}$ are full arc lengths, mixed with weight $\omega_{ab}$. Averaging over positions gives $\omega_{ab}a_{ab}+(1-\omega_{ab})b_{ab}=P_{ab}$. Set $g_{ab}^{(0)}=P_{ab}$ and $g_{ab}^{(1)}(\delta)=\mathbf1\{\delta\le P_{ab}/2\}$ at the endpoints. Edges incident to the central node remain independent after integrating over candidate positions, so~\eqref{eq:fixed_degree_composition} still applies. Conditional on connection to a type-$a$ central node at zero, draw neighbor positions independently, with each type-$h$ neighbor having density:
\begin{equation}
f_{ah}^{(\kappa)}(x\mid\mathrm{edge})=
\frac{g_{ah}^{(\kappa)}(\delta_{\mathrm{circ}}(x,0))}{P_{ah}},
\quad x\in[-1/2,1/2).
\label{eq:geometry_conditional_position}
\end{equation}
Here $\delta_{\mathrm{circ}}$ denotes shortest circular distance. For neighbor types $h_r$ and positions $\mathbf X=(X_r)_{r=1}^d$, sum conditionally independent closures into triangle degree $T_G$ and integrate to obtain the geometric reference $R_G$:
\begin{equation}
\begin{aligned}
&T_G\mid\mathbf z,\mathbf X
\overset{\mathrm d}{=}\sum_{r<s}\operatorname{Bernoulli}
\!\left(g_{h_rh_s}^{(\kappa)}(\delta_{\mathrm{circ}}(X_r,X_s))\right),\\
&R_G(t\mid d;\kappa)=\sum_a\eta_a(d)
\mathbb E_{\mathbf z,\mathbf X\mid a,d;\kappa}
[\Pr(T_G=t\mid\mathbf z,\mathbf X)].
\end{aligned}
\label{eq:geometry_conditional_reference}
\end{equation}
Shared positions induce dependence among closures after integration, without estimating actual node locations. At $\kappa=0$, this reduces to $R_\beta$.

\emph{Selecting $\kappa$.} Using the same truncation working model as Section~\ref{sec:base}, match the predicted report mean $\mu_G^{\mathrm{rep}}(\kappa)$ to the observed mean $\bar s$:
\begin{equation}
\begin{aligned}
&\mu_G^{\mathrm{rep}}(\kappa)=\frac1n\sum_v\sum_d
q_v(d)\frac{W^C_{v,d}}{w(d)}\mathbb E[T_G\mid d;\kappa],\\
&\widetilde\mu_G^{\mathrm{rep}}(\widehat\kappa)
=\operatorname{clip}_{I_G^{\mathrm{rep}}}(\bar s),
\qquad \bar s=\frac1n\sum_vs_v.
\end{aligned}
\label{eq:geometry_calibration}
% Equation label merged: eq:geometry_population_mean -> eq:geometry_calibration
\end{equation}
Solve for $\widehat\kappa\in[0,1]$. Zero-capacity terms are zero. Here $\widetilde\mu_G^{\mathrm{rep}}$ is the numerical approximation, and $I_G^{\mathrm{rep}}$ is the interval between its endpoint predictions at $0$ and $1$. Use a matching endpoint when the target equals its prediction, and $\widehat\kappa=0$ when the endpoint predictions coincide. Fix $R_G(t\mid d)=R_G(t\mid d;\widehat\kappa)$. Both references are computed analytically for $d\le2$ and by Monte Carlo integration otherwise. Neither is assumed correct across all triangle-degree ranges.


\subsection{Phase 3: Report-Driven Distribution Correction}
\label{sec:correction}
To address Challenge 2, we use the two structural references
to correct the data-driven base $R_F$ without treating
either as ground truth.
A reference may capture the relative frequencies within
a count range but misestimate its total probability,
so we adjust regional mass and within-region shape separately.
Since randomized reports provide limited evidence,
we fit a small set of correction strengths shared across
degrees using the report likelihood.

\subsubsection{Regional mass and shape}
Using shared boundaries $b_1\le b_2$, define four degree-specific regions $I_r(d)$, separating zero from positive counts:
\begin{equation}
\begin{aligned}
&I_0(d)=\{0\},\\
&I_1(d)=(0,b_1]\cap\{1,\ldots,w(d)\},\\
&I_2(d)=(b_1,b_2]\cap\{1,\ldots,w(d)\},\\
&I_3(d)=(b_2,w(d)]\cap\mathbb Z.
\end{aligned}
\label{eq:participation_regions}
\end{equation}
These are zero, body, transition, and tail. Choose $b_1,b_2$ at two prespecified quantile levels of the base population distribution $p_F(t)=\sum_d\widehat\pi(d)R_F(t\mid d)$ conditional on $t>0$, and keep them fixed during fitting. If positive mass is absent, set both to zero; empty regions are allowed. Different regions can then change without forcing the same adjustment elsewhere.

Within a region, a reference can have a useful relative shape but an inaccurate total probability. For source $s\in\{F,G,\beta\}$ and region $r$, define regional mass $m_{sr}(d)$ and normalized shape $h_{sr}(t\mid d)$:
\begin{equation}
\begin{aligned}
&m_{sr}(d)=\sum_{t\in I_r(d)}R_s(t\mid d),\\
&h_{sr}(t\mid d)=\frac{R_s(t\mid d)\mathbf1\{t\in I_r(d)\}}{m_{sr}(d)}.
\end{aligned}
\label{eq:regional_mass_shape}
\end{equation}
Set $h_{sr}=0$ when $m_{sr}=0$. For each positive region, choose initial within-region shape $h_{0r}$ as the first nonempty shape in the order $F,G,\beta$, or zero if none exists. Set $h_{00}(t\mid d)=\mathbf1\{t=0\}$. Transferring mass along $h_{0r}$ avoids changing shape. %at the same time.

\subsubsection{Shared correction parameters}
Suppressing $(d,t)$, define total positive mass $m_s^+=\sum_{r=1}^3m_{sr}$ and its regional proportions $\phi_{sr}=m_{sr}/m_s^+$. Set $\phi_{Fr}=0$ if $m_F^+=0$; a reference with no positive mass uses $\phi_{Fr}$. Three changes $\Delta_j(t\mid d)$ to count probabilities transfer mass between zero and positive counts and redistribute it across positive regions:
\begin{equation}
\begin{aligned}
&\Delta_1=(m_G^+-m_F^+)\left(\sum_{r=1}^3\phi_{Gr}h_{0r}-h_{00}\right),\\
&\Delta_2=m_F^+\sum_{r=1}^3(\phi_{Gr}-\phi_{Fr})h_{0r},\\
&\Delta_3=m_F^+\sum_{r=1}^3(\phi_{\beta r}-\phi_{Fr})h_{0r}.
\end{aligned}
\label{eq:structural_directions}
\end{equation}
Only $\Delta_1$ changes total positive mass. For the four within-region shape directions, we fix the assignment: geometric in the body, both in the transition, and $\beta$ in the tail.
\begin{equation}
\begin{alignedat}{2}
&\Delta_4=m_{G1}(h_{G1}-h_{01}), &\quad& \Delta_5=m_{G2}(h_{G2}-h_{02}),\\
&\Delta_6=m_{\beta2}(h_{\beta2}-h_{02}), &\quad& \Delta_7=m_{\beta3}(h_{\beta3}-h_{03}).
\end{alignedat}
\label{eq:structural_shape_directions}
\end{equation}
Each direction sums to zero; shape directions also conserve regional mass. The fixed shape allocation does not restrict the mass directions. Set $R_{\mathbf x}=R_F+\sum_{j=1}^7x_j\Delta_j$, where $\mathbf x=(x_1,\ldots,x_7)$ contains seven strengths shared across degrees. They scale changes, not mixture components, and need not sum to one.

Bounds $0\le x_j\le1$ alone do not ensure nonnegativity. For each $r=0,\ldots,3$, order the three fixed shapes as $(h_{0r},h_{Gr},h_{\beta r})$ to form $\{h_{d\ell}\}_{\ell=1}^{12}$. Expanding the displayed directions in this order gives the following representation and feasible set $\mathcal C$:
\begin{equation}
\begin{aligned}
&R_{\mathbf x}(t\mid d)=\sum_{\ell=1}^{12}
[\mathbf c_d+\mathbf B_d\mathbf x]_\ell h_{d\ell}(t),\\
&\mathcal C=\{\mathbf x\in[0,1]^7:
\mathbf c_d+\mathbf B_d\mathbf x\ge0\ \forall d\in\mathcal D\}.
\end{aligned}
\label{eq:joint_feasible_domain}
\end{equation}
Here $\mathbf c_d\in\mathbb R^{12}$ has regional blocks $(m_{Fr},0,0)$, and column $j$ of $\mathbf B_d\in\mathbb R^{12\times7}$ contains the coefficients of $\Delta_j$ in these fixed shapes. These twelve coefficients express constraints; only seven parameters are learned. The constraints and zero-sum directions preserve nonnegativity, normalization, and the degree-specific count range; $\mathbf x=0$ retains the base.

\subsubsection{Learning from the observed reports}
For $w(d)>0$, define the report likelihood $\mathcal K_v(Q_v\mid d,t)=\ell_v(d,t/w(d))$; for $w(d)=0$, set $\mathcal K_v(Q_v\mid d,0)=\ell_v(d,0)$. Keep the base fit's truncation model and degree support. With report likelihood $f_v(\mathbf x)$, fit:
\begin{equation}
\begin{aligned}
&f_v(\mathbf x)=\sum_dq_v(d)\sum_{t=0}^{w(d)}
R_{\mathbf x}(t\mid d)\mathcal K_v(Q_v\mid d,t),\\
&\mathbf x^\star\in\arg\min_{\mathbf x\in\mathcal C}
\left\{-\sum_v\log f_v(\mathbf x)+\frac12\|\mathbf x\|_2^2\right\}.
\end{aligned}
\label{eq:fusion_optimization}
\end{equation}
This composite likelihood selects corrections from reports; the penalty shrinks them toward the base. The base, references, boundaries, degree probabilities, and report likelihoods are fixed during fitting. Affine report densities and linear constraints give a convex objective on positive densities. The corrected conditional distribution is $R_J=R_{\mathbf x^\star}$.


\subsection{Phase 4: Posterior Aggregation and Final Release}
Although $R_J$ provides a shared degree-conditioned distribution, individual reports may favor different triangle-degree values. 
Phase 4 uses these reports for posterior refinement and aggregation to produce the final TDD.

\label{sec:aggregation}
Using $\pi(d)$ and $R(t\mid d)$ defined in Section~\ref{sec:pre}, the population histogram satisfies:
\begin{equation}
H(t)=n\sum_{d=0}^{D_{\max}}\pi(d)R(t\mid d).
\label{eq:problem-conditional-factorization}
\end{equation}
Conditioning on degree preserves the dependence between wedge count and closure rate. The factor $n$ returns node counts rather than normalized probabilities.

Regional correction links count probabilities through seven directions, but each report can favor individual counts differently. Freeze $R_J$ and compute joint node-degree and triangle-degree probabilities $r_v^J(d,t)$ by one posterior reweighting~\cite{efron2014two}:
\begin{equation}
r_v^J(d,t)=\frac{q_v(d)R_J(t\mid d)\mathcal K_v(Q_v\mid d,t)}
{f_v(\mathbf x^\star)}.
\label{eq:frozen_joint_posterior}
\end{equation}
To preserve the estimated degree proportions $\widehat\pi$, normalize within each degree to obtain $R^\star(t\mid d)$, then aggregate:
\begin{equation}
\begin{aligned}
&R^\star(t\mid d)=\frac{\sum_vr_v^J(d,t)}
{\sum_v\sum_{t'=0}^{w(d)}r_v^J(d,t')},\\
&\widehat H_{\mathrm{raw}}(t)=n\sum_d\widehat\pi(d)R^\star(t\mid d).
\end{aligned}
\label{eq:final_gdd}
\end{equation}
If the within-degree normalization denominator for $R^\star(t\mid d)$ is zero, set $R^\star(t\mid d)=R_J(t\mid d)$. This single update can move beyond the finite correction family, but cannot add mass where $R_J=0$. The candidate estimate $\widehat H_{\mathrm{raw}}$ has total count $n$.

For the transport safeguard, let $F_{\mathrm{dir}}$ be the empirical cumulative distribution function (CDF) of $\mathsf b_{W_{\max}}(s_v)$ over the original relational reports $s_v$, using the clipping and rounding rule in Eq.~\eqref{eq:direct-gdd}. The estimated wedge-count CDF is $F_{\mathrm{cap}}(t)=\sum_d\widehat\pi(d)\mathbf1\{w(d)\le t\}$. The reference histogram has counts $H_A(t)=n[F_A(t)-F_A(t-1)]$, where $F_A(-1)=0$ and
\begin{equation}
F_A(t)=\max\{F_{\mathrm{dir}}(t),F_{\mathrm{cap}}(t)\}.
\label{eq:pg-reference-cdf}
\end{equation}
Let $\Delta=W_1(\widehat H_{\mathrm{raw}}/n,H_A/n)$, where $W_1$ is the first Wasserstein distance with ground cost $|t-u|$ between triangle degrees. For public $\tau>0$, set $\alpha=1$ if $\Delta\le\tau D_{\max}$ and $\alpha=\tau D_{\max}/\Delta$ otherwise. PrivTDD releases:
\begin{equation}
\widehat H=\alpha\widehat H_{\mathrm{raw}}+(1-\alpha)H_A.
\label{eq:pg-guard}
\end{equation}
This post-processing uses only released reports and public parameters, so it incurs no additional privacy loss.
